\documentclass[12pt,a4paper]{article}
\usepackage[utf8]{inputenc}
\usepackage[greek]{babel} % for greek
%\usepackage[T1]{fontenc}
\usepackage{amsmath}
\usepackage[svgnames]{xcolor}
\usepackage{amsfonts}
\usepackage{amssymb}
\usepackage{graphicx}
\usepackage{hyperref}
\usepackage[left=1.50cm, right=1.50cm, top=3.00cm, bottom=3.00cm]{geometry}
%\author{Stratos Pateloudis}
\title{Μια πρώτη εκτίμηση για τις ελιές}
\begin{document}
\maketitle
\section{Ανάλυση}
Η ανάλυση βασίζεται σε δεδομένα που συλλέχθηκαν απο το 2004 εως και σήμερα. Η πιστώτητα των δεδομένων είναι μεγαλύτερη για τα κιλά των ελιών και του λαδιού στα οποία και θα επικεντρωθούμε. Aς ξεκινήσουμε παρουσιάζοντας ένα γράφημα για την εξέλιξη σε αυτά τα χρόνια όπως φαίνεται στο Σχήμα \ref{fig:general}. Απο το σχήμα παρατηρείται μια αυξομείωση η οποία αναμένεται, δηλαδή κάθε δεύτερο χρόνο οι ελιές είναι σε μεγάλη καρποφορία ενώ στον ενδοιάμεσο χρόνο η παραγωγή είναι πολύ μικρή (φαινόμενο παρενιαυτοφορίας). Όπως θα δούμε όμως υπάρχει μια μέθοδος να έχουμε ελιές κάθε χρόνο, υψηλής και σταθερής παραγωγικότητας. 
\begin{figure}[ht!]
	\centering 
	\includegraphics[scale=0.7]{general_analysis}
	\caption{Ο άξονας $y$ είναι σε απόλυτους αριθμούς και φαίνεται η εξέλιξη των ελιών, του λαδιού, της απόδοσης και της τιμής κιλού. Η απόδοση και η τιμή κιλού είναι πολλαπλασιασμένη επί 1000 για αισθητικούς λόγους. Το έτος 2007 λάδι, ελιές και απόδοση έχουν μηδενιστεί λόγω έλειψης δεδομένων.}
	\label{fig:general}
\end{figure}

 Απο αυτήν την πρώτη εκτίμηση μπορούμε να κάνουμε μια περαιτέρω ανάλυση για να δούμε συσχετίσεις με τις διάφορες μεταβλητές. Αυτό φαίνεται στο Σχήμα \ref{fig:corr}. Παρατηρείται (όπως άλλωστε αναμένεται) μια γραμμική, δηλαδή ισχυρή συσχέτιση μεταξύ των κιλών των ελιών και του λαδιού που προτείνει μια σταθερή απόδοση όπως θα δούμε. 
\begin{figure}[ht!]
	\centering 
	\includegraphics[scale=0.8]{correl_olives}
	\caption{Ένας πίνακας που δίνει την συσχέτιση δεδομένων για τις διάφορες μεταβλητές, όπως Ελιές, Λάδι, Απόδοση, Τιμή κιλού. Γραμμική αύξηση σημαίνει ισχυρή συσχέτιση, όπως πχ Ελιές και Λάδι. Η διαγώνιος δείχνει που κατανέμεται κάθε μεταβλητή, πχ οι ελιές κατανέμονται στα 1000 κιλά ενώ το λάδι περί τα 250 κιλά.}
	\label{fig:corr}
\end{figure}
Μπορούμε να προχωρήσουμε περισσότερο υποθέτωντας ότι υπάρχουν περίπου 185 παραγωγικά δέντρα και στα δυο χωράφια (134 στο πάνω και 51 στο κάτω). Τότε μπορούμε να βρούμε αρχικά τον μέσο όρο των ελιών και του λαδιού ανα χρόνο. Αυτά δίνονται στον πίνακα \ref{tab:means}.
\begin{table}[ht!]
	\centering
	\begin{tabular}{|c|c|}
		\hline 
		Μέσος όρος ελιών	& 1484.7 κιλά \\ 
		\hline 
		Μέσος όρος Ελαιόλαδου	&  348.5 κιλά\\ 
		\hline
		Μέση απόδοση & 0.23 $\left(\frac{1}{4.25}\right)$\\
		\hline 
		Ελιές / δέντρο & 8.13 κιλά \\
		\hline
		Μέσος όρος λαδιού / δέντρο & 1.91 κιλά\\
		\hline
	\end{tabular} 
 \caption{Γενικοί μέσοι όροι απο όλα τα χρόνια. Απο τον μέσο όρο εξαιρείται η χρονιά του 2007 που δεν υπάρχουν στοιχεία.} \label{tab:means}
\end{table}
H χρονιά 2007 δεν λαμβάνεται υπόψιν στο μέσο όρο. Από τον πίνακα διαπιστώνουμε ότι η απόδοση ανά δέντρο είναι κάτω του μέσου όρο και αυτό μπορεί πιθανώς να εξηγηθεί από το έλλειμμα στοιχείων (οργανικών και ανόργανων) που υπάρχουν στα χωράφια και αφαιρούνται κάθε καλή χρονιά που η παραγωγή είναι καλή. Θα πρέπει να επανατροφοδοτούμε τα δέντρα συνεχώς και όταν το χρειάζονται.

Ένας άλλος τρόπος να βρούμε την μέση απόδοση είναι να δημιουργήσουμε μια γραμμική συνάρτηση που συσχετίζει τις ελιές και το λάδι. Αναμένουμε να μας δώσει έναν λόγο κοντά στο $1/4$. Η συνάρτηση που κατασκευάζουμε είναι 
\begin{align*} 
y = a\cdot x+ b,
\end{align*} 
όπου $y$ είναι το λάδι σε κιλά, $a$ είναι ο λόγος, $x$ είναι οι ελιές σε κιλά  και $b$ είναι μια σταθερά. Το αποτέλεσμα φαίνεται στο Σχήμα \ref{fig:fit}.

\begin{figure}[ht!]
	\centering 
	\includegraphics[scale=0.6]{fit_oil_olives}
	\caption{Στατιστικό μοντέλο που περιγράφει την εξάρτηση των ελιών και του λαδιού, δίνοντας τον λογο. Οι κουκίδες είναι τα δεδομένα ενώ η κόκκινη γραμμή είναι η συνάρτηση $y=a\cdot x +b$. }
	\label{fig:fit}
\end{figure}
Από το σχήμα συμπαιρένουμε ότι η γραμμική συνάρτηση που περιγράφει την συσχέτιση του λαδιού και των ελιών δίνεται ως 
\begin{align*}
y=0.23\cdot x +6.18,
\end{align*}
όπου $y$ είναι το λάδι σε κιλά και $x$ είναι οι ελιές σε κιλά δίνοντας ένα λόγο $0.23$ όπως αναμενόταν που αντιστοιχεί περίπου σε $1/4.34$. Αυτή είναι μια χονδροειδής προσέγγιση όμως ανταποκρίνεται στην πραγματικότητα αφού ο μέσος όρος προσεγγίζει τον πραγματικό (πχ Πίνακας \ref{tab:means}). Σε περίπτωση που αλλάξει το μοντέλο στο μέλλον σημαίνει ότι η απόδοση έχει επηρεαστεί. 

\section{Συμπέρασμα}
Σύμφωνα με την βιβλιογραφία \textbf{τα επίπεδα του αζώτου συγκεκριμένα παίζουν μεγάλο ρόλο στην συνεχή και αυξημένη παραγωγή}, που σημαίνει οτι πρέπει να επικεντρωθούμε σε αυτό. Σύμφωνα με τελευταία στοιχεία, εμπλουτισμός του χωραφιού με αυξημένα επίπεδα αζώτου \footnote{Μια άλλη έρευνα προτείνει 15 κιλά/στρέμμα άζωτο ως ιδανικά κιλά απο την άποψη κόστους/απόδοσης.} (περίπου 12 κιλά/στρέμμα καθαρού αζώτου συνολικά) φέρνει καλύτερες αποδόσεις στην παραγωγή. Ίσως όχι άμεσα αλλά σε βάθος 2-4 χρόνων σίγουρα.  

Παράλληλα ο \textbf{ετήσιος εμπλουτισμός με βόριο} κρίνεται σημαντικός, ίσως όχι τόσο στην παραγωγή αλλά στην ενδυνάμωση του δέντρου. 

Μια πρόσφατη μελέτη (2021) προτείνει οτι η υπερενίσχυση του χωραφιού με άζωτο απέφερε την μέγιστη καρποφορία ανά ελιά, ξεκινώντας γύρω στα 10 κιλά ανά ελιά (όπως και με εμάς, βλέπε Πίνακα \ref{tab:means}) την πρώτη χρονιά σε όλες τις μεθόδους (2017) και φτάνοντας τα 20+ κιλά ανά δέντρο το 2020 για την μέθοδο υψηλών επιπέδων αζώτου (Σχήμα \ref{fig:results_high_N}). H χρήση του αζώτου κρίνεται ιδιαίτερα σημαντική για μια σχετικά σταθερή και συνεχώς αυξανόμενη παραγωγή κιλών καρπού \textbf{κάθε χρόνο}. Το σημαντικό εδώ είναι ότι κάθε χρόνο οι ελιές καρποφορούσαν με οποιαδήποτε μέθοδο, αν και τις κακές χρονιές (χρονιές παρενιαυτοφορίας) στις υπόλοιπες μεθόδους υπήρχε μια μικρή έως μηδαμινή μείωση της παραγωγής.
\begin{figure}[ht!]
	\centering
  \includegraphics[scale=0.6]{yield}
  \caption{Ετήσια παραγωγή ανα δέντρο σε κιλά καρπού από την έρευνα του 2021. Οριζόντια είναι οι 5 διαφορετικές μέθοδοι που εφαρμόστηκαν όπου η τελευταία (δεξιά στο πλαίσιο) μέθοδος είναι η υπερεμπλούτιση με άζωτο. Ενώ όλες οι μέθοδοι ξεκίνησαν με περίπου 10 κιλά ανα ελιά (όπως και στην δική μας περίπτωση, Πίνακας \ref{tab:means}) πολύ γρήγορα (ήδη απο την 2η χρονιά) η εφαρμογή της τεχνικής υψηλών επιπέδων αζώτου ξέφυγε απο τις άλλες, φτάνοντας το 2020 περίπου σε 70 κιλά/δέντρο \textbf{συνολικά και τα τέσσερα χρόνια}. Οι υπόλοιπες μέθοδοι χρησιμοποίησαν οργανικά κατάλοιπα, πχ κοπριά καθώς και στάχτη σε ποσότητες ώστε να περιέχουν 5 κιλά/στρέμμα άζωτο. Η μέθοδος υψηλού αζώτου φαίνεται να δίνει χονδρικά διπλάσια και σταθερή ετήσια παραγωγή απο τις άλλες μεθόδους. }
  \label{fig:results_high_N}
\end{figure}


Η έρευνα έγινε σε δέκα στρέμματα ξερικών ελιών (22 χρονών) έχοντας περίπου 204  δέντρα ($7\times 7$ μέτρα αποστάσεις). H μέση ετήσια θερμοκρασία της περιοχής ήταν 14 βαθμοί, με μέση ετήσια βροχόπτωση 509$mm$ ενώ το χωράφι βρισκόταν σε ύψος 240 μέτρα απο την επιφάνεια της θάλασσας. 

Αυτά τα χαρακτηριστικά βρίσκονται κοντά σε αυτά της Ζίχνης, όπου έχει μέση θερμοκρασία 15.84 βαθμούς Κελσίου, περίπου 728$mm$ ετήσια βροχόπτωση\footnote{Στοιχεία από \href{https://www.worldweatheronline.com/nea-zichni-weather-averages/central-macedonia/gr.aspx}{\color{red} εδώ}.}. Παραδοσιακά η περίοδος των βροχοπτώσεων είναι στο τέταρτο τρίμηνο του έτους με
σημαντική εξαίρεση την χρονιά του 2023. Επίσης το ύψος (στο χωράφι) είναι περίπου 360μ.  

Η υπερεμπλούτιση του χωραφιού στην έρευνα έγινε χρησιμοποιώντας αρχικά βασική λίπανση, απο άζωτο, φώσφορο και κάλιο με εφαρμογή 5 κιλά/στρέμμα  (ή 245 γραμμάρια έκαστο/δέντρο) καθώς επίσης και με επιπλέον 7 κιλά/στρέμμα (ή 340 γραμμάρια/δέντρο) καθαρό άζωτο. Επιπλέον, 0.2 κιλά/στρέμμα βόριο (ή 9.8 γρ/δέντρο), και όλα αυτά εφαρμόστηκαν ετησίως πριν τον Μάρτιο.

Για την βασική λίπανση χρησιμοποιήθηκαν λιπάσματα με αναλογία $10\%/10\%/10\%$ σε άζωτο, φώσφορο και κάλιο ενώ η υπερεμπλούτιση με το άζωτο έγινε μέσω νιτρικής αμμωνίας με περιεκτικότητα 20.5$\%$ σε άζωτο.

Απο τα προηγούμενα, φαινεται, οτι αυτό που πρέπει να κάνουμε εμείς είναι να ρίξουμε βασικό λίπασμα, τόσο ώστε να αντιστοιχεί σε περίπου 250 γρ./δέντρο σε καθαρό άζωτο, φώσφορο και κάλιο. Επίσης εφαρμογή επιπλέον 350 γρ./δέντρο μόνο άζωτο μέσω πχ νιτρικής αμμωνίας, καθώς επίσης και 10 γρ./δέντρο καθαρό βόριο. Το πως εφαρμόζεται αυτή η φόρμουλα δεν είναι ξεκάθαρο, αρκεί να επιτυγχάνονται αυτές οι ποσότητες ανα ρίζα ετησίως. Ωστόσο, σε ξερικές ελιές μπορεί να γίνει σε δύο δόσεις, πχ τον χειμώνα εφαρμογή μόνο του βασικού μαζί με βόριο για την εκμετάλλευση των βροχών, και την άνοιξη εφαρμογή της νιτρικής αμμωνίας.

Συλλέγοντας όλα τα δεδομένα έχουμε οτι \textbf{για ένα δέντρο ετησίως} χρειαζόμαστε την φόρμουλα στον Πίνακα \ref{tab:formula}.
\begin{table}[ht!]
	\centering
	\begin{tabular}{|c|c|}
		\hline
		\textbf{Στοιχείο} & \textbf{Στη ρίζα}\\
		\hline 
		\hline
		Άζωτο &  600 γραμμάρια \\ 
		\hline 
		Φώσφορο &  250 γραμμάρια \\ 
			\hline 
		Κάλιο &  250 γραμμάρια \\ 
		\hline 
		Βόριο &  10 γραμμάρια\\ 
		\hline 
	\end{tabular} 
\caption{Φόρμουλα που πρέπει να εφαρμοστεί ετησίως ανά δέντρο. Τα ποσά αναφέρονται σε καθαρές ποσότητες στοιχείων τα οποία παρέχονται απο τα λιπάσματα με την αντίστοιχη αναλογία.}\label{tab:formula}
\end{table}

Σύμφωνα πάλι με την διεθνή βιβλιογραφία φαίνεται πως τα επίπεδα αζώτου στο δέντρο (με ψεκασμό) δεν επηρεάζουν την ποιότητα του λαδιού. 

\subsection{Άζωτο}
Το άζωτο φαίνεται να είναι το σημαντικότερο στοιχείο για κηπουρικές καλιέργιες.  Σχετικά πρόσφατα διαπιστώθηκε ότι οι ελιές ανταποκρίνονται θετικά στην παροχή αζώτου τόσο οι ξερικές όσο και οι ποτιστικές. Οι ξερικές ωστόσο έχουν τα μεγαλύτερα οφέλη καθώς η παροχή αζώτου φαίνεται να τις κάνει πιο ανθεκτικές στην λειψυδρία. Απο την άλλη είναι δύσκολο να επεμβαίνουμε και να διαχειριζόμαστε τα επίπεδα αζώτου στο χώμα, στο ριζικό σύστημα και κατ' επέκταση μέσω της απορρόφησης στα φύλλα, σε αντίθεση με τις ποτιστικές που μπορεί το άζωτο να εισαχθεί πιο εύκολα στο ριζικό σύστημα μέσω του ποτίσματος. Για αυτό κρίνεται απαραίτητο στις ξερικές ελιές να εφαρμοστεί κατα την διάρκεια των βροχοπτώσεων ώστε να εκμεταλευτούμε το νερό για την διάσπαση και μεταφορά (ή απορρόφηση) του αζώτου στο ριζικό σύστημα. 

Τα επίπεδα αζώτου εν γένει θα πρέπει να συμμορφώνονται σε κάποιους κανόνες ώστε να μην έχουμε έλλειψη αλλά και να μην δημιουργούμε τοξικές συνθήκες για την ελιά. Γενικώς το εύρος μεταξύ των δύο αυτών ορίων είναι αρκετά μεγάλο αλλά σαν γενικό κανόνα μπορούμε να ακολουθήσουμε την βιβλιογραφία εφαρμόζοντας 14-15 κιλά συνολικού, καθαρού αζώτου ανα στρέμμα. Αυτό αντιστοιχεί σε 680-735 γραμμάρια ανά ρίζα καθαρό άζωτο.  Το πολύ άζωτο, σε ποσοστό πχ μεγαλύτερο των 30 κιλών ανα στρέμμα επηρρεάζει αρνητικά τόσο την ποιότητα του λαδιού όσο και την παραγωγή, ενώ από την άλλη επηρρεάζει θετικά την παραγωγή νεων βλαστών. Μια τέτοια μέθοδος (δηλαδή πολύ άζωτο) μπορεί να συνεισφέρει θετικά στην ανάπτυξη νεαρών ελαιόδεντρων. Μια συνολική εικόνα των προηγουμένων δίνεται στο Σχήμα \ref{fig:N_effect}.
\begin{figure}[ht!]
\centering 
\includegraphics[scale=0.2]{N_effect}
\caption{Σχηματική αναπαράσταση της συνεισφοράς του αζώτου σε διάφορες μεταβλητές, όπως την παραγωγή και ποιότητα λαδιού. H ποιότητα λαδιού δεν αφορά μόνο την οξύτητα αλλά και τα διάφορα στοιχεία που περιέχει.}
\label{fig:N_effect}
\end{figure}  
Σημαντικό είναι επίσης, το γεγονός ότι με υπερβολική χρήση αζώτου αυτό μπορεί να συγκεντρωθεί κάτω απο το ριζικό σύστημα και να μην απορροφηθεί απο το δέντρο.

Ένας άλλος αρνητικός παράγοντας της υπερβολικής χρήσης αζώτου είναι επίσης η περιβαλλοντική μόλυνση ως oξείδιο του αζώτου.

\subsection{Φώσφορος}
Μέχρι πρόσφατα δεν υπήρχαν επαρκή στοιχεία για την θετική ή αρνητική συνεισφορά του φωσφόρου στην παραγωγή της ελιάς. Ωστόσο, με έρευνες της τελευταίας δεκαετίας παρατηρήθηκε η θετική συνεισφορά του φωσφόρου στην παραγωγή της ελιάς. Συγκεκριμένα η εμπλούτιση του εδάφους με φώσφορο συνεισφέρει επίσης στην γρήγορη και συνεχή απορρόφηση του φωσφόρου απο το δέντρο καθώς το φώσφορο εν γένει κινείται και απορροφάται σχετικά δύσκολα απο το δέντρο. Οπότε η συνεχής παρουσία φυσιολογικών επιπέδων φωσφόρου στο ριζικό σύστημα κρίνεται σημαντική. 

Λόγω της δυσκολίας της μεταφοράς του φωσφόρου μέσα και ανάμεσα στο ριζικό σύστημα του δέντρου, και λαμβάνοντας υπόψιν έρευνες που παρουσιάζουν θετική συχέτιση μεταξύ νερού και κίνηση φωσφόρου στο ριζικό σύστημα, κρίνεται επίσης αναγκαίο η εφαρμογή του φωσφόρου κατά την περίοδο βροχοπτώσεων. 

\subsection{Κάλιο}
Το κάλιο φαίνεται να είναι ένα απο τα πιο σημαντικά στοιχεία για την ελιά (!) ίσως λόγω της συγκέντρωσης του στον καρπό της ελιάς (στο φρούτο και όχι στο κουκούτσι), και επίσης επειδή σε ξερικές ελιές παρατηρείται η έλλειψη καλίου πιο συχνά. Η έλλειψη καλίου μεταξύ άλλων εμποδίζει την φωτοσυνθετική ικανότητα του δέντρου, η οποία ενισχύεται και από το άζωτο, που σημαίνει ότι το κάλιο πρέπει να εφαρμόζεται παράλληλα με το άζωτο ώστε να βοηθάει στην φωτοσύνθεση και επακόλουθα στην παραγωγή καρπού. Υπερβολικά επίπεδα καλίου φαίνεται ότι δεν δημιουργούν πρόβλημα, ούτε στο χωράφι ούτε στην ατμόσφαιρα.     


\subsection{Διαφυλλική λίπανση-πρόγραμμα έτους}
Σύμφωνα με την ανάλυση χώματος και φύλλων πάμε αποκλειστικά σε διαφυλλικούς ψεκασμούς θρέψης, ενώ η χρήση ατταπουλγίτη κρίνεται απαραίτητη. 

Οι ψεκασμοί που κάνουμε περιέχουν τα παραπάνω στοιχεία $(Ν, P, K, B)$ ενώ η εφαρμογή τους γίνεται σε διαφορετικές περιόδους.


\begin{itemize}
	\item  Διαφυλλικοί ψεκασμοί με άζωτο (Ν) προ-ανθικά, ανθικά, μετα την έναρξη καρπόδεσης καθώς και μετά την σκλήρυνση του πυρήνα, φαίνεται πως ευνοεί το δέντρο, επιταχύνοντας την ανάπτυξη βλάστησης και καλυτερεύουν την ποσότητα καρπών. Ιδιαίτερα σε ξερικές ελιές, ο ψεκασμός σε συνθήκες στρές (όπως το καλοκαίρι) ενισχύει την ανάπτυξη του δέντρου, την ποιότητα του ελαιολάδου, και την διαθεσιμότητα του αζώτου στις κρίσιμες αυτές περιόδους. Ψεκασμοί γίνονται συνήθως με αραίωση 1-4\% και χρησιμοποιώντας ουρία (καθώς έχει μεγαλύτερη συγκέντρωση αζώτου).
	\item Παρόμοια, ο φώσφορος ψεκάζεται με σκευάσματα σε αραίωση 1-3\% ενώ το κάλιο πάλι σε παρόμοια αραίωση (πχ 1-3\% $KNO_3$) κατά την δεύτερη και τρίτη φάση ωρίμανσης του καρπού (στο μεγάλωμα των ελιών και στην έναρξη αλλαγής χρώματος, περίπου Αύγουστο σε εμάς) είχε θετική επίδραση στο μέγεθος των καρπών και στην ελαιοπεριεκτηκότητα των καρπών. Η απορρόφηση του καλίου γίνεται κυρίως μέσα απο τα νεαρά φύλλα. 
	\item  Το μαγνήσιο έχει επίσης θετική επίδραση σε αναλογία 1.5\% σε κορονέικη, βελτιώνοντας την βλάστηση, ανάπτυξη καρπών, και ποιότητα ελαιολάδου.
	\item Το βόριο ίσως είναι το πιο σημαντικό. Σε ελιές χαλκιδικής \underline{με τροφοπενία βορίου} (για επάρκεια βορίου λιγότερο προφανώς) ψεκασμοί προ-ανθικά και μετα-ανθικά, (15 μέρες πριν και 15 μετά) έδειξαν θετική επίδραση στην ποσότητα καρποσυγκράτησης και ανάπτυξης ελιών. Συγκεκριμένα χρησιμοποιήθηκαν ψεκασμοί με σκεύασμα Β σε αναλογία $567mg$ καθαρού στοιχείου $B/L$. Με μια μέση περιεκτικότητα $10\%$ αυτό αντιστοιχεί σε 5.67 κιλά λιπάσματος ανά 1000$L$.
 
\end{itemize}

\section{Εφαρμογή φόρμουλας}
\subsection{2022}
Σκάψιμο περιμετρικά των ριζών στο ύψος των φύλλων σε 185 δέντρα συνολικά (134 στο πάνω χωράφι και 51 στο κάτω) και τοποθέτηση του λιπάσματος μέσα στο αυλάκι. Στο πάνω χωράφι στα πρώτα δέντρα της δεύτερης αράδας εκ της νέας εισόδου η τοποθέτηση έγινε στο εσωτερικό του αυλακιού. 

\begin{itemize}
	\item Εφαρμογή βασικού τα Χριστούγεννα με αναλογία $20/8/5 \%$ σε άζωτο φώσφορο και κάλιο αντίστοιχα με 2κλ/δέντρο. Εφαρμογή Βόρακα $10\%$ με 100γρ/δέντρο και εφαρμογή θειικού καλίου με περιεκτικότητα $50\%$ σε $Κ_2Ο$\footnote{$\%K_2O\times 0.83=\%K$} με 350γρ/δέντρο. 
	\item Κλάδεμα μέσα Μαρτίου με εφαρμογή νιτρικής αμμωνίας (33.4\% Ν) περίπου 600γρ/δέντρο στο αυλάκι και σκεπασμα.  
	\end{itemize}

\begin{table}[ht!]
	\centering
	\begin{tabular}{|c|c|}
		\hline
		\textbf{Στοιχείο} & \textbf{Στη ρίζα}\\
		\hline 
		\hline
		Άζωτο &  600 γρ \\ 
		\hline 
		Φώσφορο &  140 γρ \\ 
		\hline 
		Κάλιο &  245 γρ \\ 
		\hline 
		Βόριο &  10 γρ\\ 
		\hline 
	\end{tabular} 
	\caption{Φόρμουλα που εφαρμόστηκε στις ρίζες την περίοδο 2022.}\label{tab:formula}
\end{table}
Στο πάνω χωράφι στις 2μιση αράδες από πάνω εφαρμόστηκε παραπάνω κάλιο, περίπου 20 κιλά ($Κ_2Ο$) συνολικά σε αυτά τα δέντρα, λόγω του περισσεύματος από τα λιπάσματα των Χριστουγέννων. Περίπου 220γρ/δέντρο επιπλέον υποθέτωντας ότι είναι 75 δέντρα.


\begin{itemize}
	\item Μετά την συγκομιδή του 22-23, ασχολούμαστε με το πάνω χωράφι μόνο. 134 δέντρα πριν, αλλά έγινε φύτευση 12 διετών δέντρων ποικιλίας πετροελιάς στις 19 Απριλίου 2023, 146 συνολικά.
	\item Η φύτευση έγινε με 800γρ ατταπουλγίτη και περίπου 500-800 γρ βασικού στο λάκκο φύτευσης και ακολούθησαν βροχές. 
	%\item  Πολλές βροχές όλο τον Μάιο.
\end{itemize}

\subsection{2023}
Εφαρμογή μόνο στο πάνω χωράφι, περίπου σε 130 παραγωγικά δέντρα. Ατταπουλγίτης, στις 2 πάνω αράδες 4 κιλά ανα δέντρο, και προς τα κάτω 3,3,2,2 και σκόρπισμα μεχρι την κόμη του δέντρου με μέσο όρο περίπου 3200 γρ/δέντρο. Βασικό λίπασμα με αναλογία 20/10/0 με δυο κιλά ανα δέντρο, και  1 κιλό/δέντρο καλιομαγνήσιο υπο της μορφής 30/10\% σε $K_2O$ και $MgO$\footnote{$\%MgO\times 0.603 = \%Mg$} αντίστοιχα. Επίσης, εφαρμογή $45$ γρ. βόρακα ανα ρίζα. Η εφαρμογή των λιπασμάτων έγινε με μίξη σε κουβά, άδειασμα στο αυλάκι του προηγούμενου έτους και σκέπασμα. 
\begin{table}[ht!]
	\centering
	\begin{tabular}{|c|c|}
		\hline
		\textbf{Στοιχείο} & \textbf{Στη ρίζα (γρ)}\\
		\hline 
		\hline
		Ατταπουλγίτης &  3200 \\ 
		\hline
		Άζωτο &  600 \\ 
		\hline 
		Φώσφορο &  200 \\ 
		\hline 
		Κάλιο &  249 \\ 
		\hline 
		Μαγνήσιο &  60.3 \\ 
		\hline 
		Βόριο &  4.5 \\ 
		\hline 
	\end{tabular} 
	\caption{Φόρμουλα που εφαρμόστηκε στις ρίζες την περίοδο 2023.}
\end{table}  
Επισημαίνεται ότι και τα δύο λιπάσματα βασικό και καλιομαγνήσιο περιείχαν θείο σε ποσοστό 23\% και 42.5\% αντίστοιχα.  

Η εφαρμογή του αζώτου μετά το κλάδεμα έγινε με ουρία (στο αυλάκι)  και τον απρίλιο, στην νέα βλάστηση ψεκάστηκαν με $20/20/20$ σκεύασμα διαφυλλικά (100$ml$/100$L$, μισή ελάχιστη συνιστώμενη δοσολογία). 

Αρχές Ιουλίου ψεκασμός με $NPK ~20\%$ και δοσολογία $150ml/100L$, ενώ βόριο (11\%)  $50ml/100L$. Αυτές οι δοσολογίες είναι κάτω απ τις συνιστώμενες από τον κατασκευαστή ($200-400 ml / 100 L, 150-180 ml / 100 L$ αντίστοιχα).
 

\end{document}